\chapter{转换实验设计与结果分析}
\label{chap:6}
\chapterguide{本章为待开展或待统一复现的实验框架。先定义样例、预期、指标与失败口径，再填入实际数据；历史记录不能替代本轮实验结果。}

\section{实验环境、路线与测试样例}
\label{sec:6-1}

\subsection{软件版本、代码版本与容差配置}
\label{subsec:6-1-1}
\outlinehint{固定三套CAD、两仓配套代码、模板、单位、操作系统和比较工具；记录实际配置，不使用未核实的版本号和容差。}

\subsection{同系统基线与跨系统转换路线}
\label{subsec:6-1-2}
\outlinehint{设置三条同系统读取重建基线和六条有向跨系统路线；“往返转换”仅用于实际执行源到目标再到源的试验。}

\subsection{五类特征样例与组合模型}
\label{subsec:6-1-3}
\outlinehint{从现有31组建议案例中筛选可复现样例，安排单因素、参数组合、连续依赖与预期失败模型，先登记预期。}

\section{评价方法与统计指标}
\label{sec:6-2}

\subsection{特征结构、参数与约束检查}
\label{subsec:6-2-1}
\outlinehint{检查特征类型、依赖和参数含义，容许有解释的辅助节点；将尺寸是否驱动、引用是否关联作为独立项目。}

\subsection{曲线、平面与实体的几何比较}
\label{subsec:6-2-2}
\outlinehint{草图检查曲线与位置，基准面检查位置方向，封闭实体检查整体量和双向材料差；处理拓扑分段与接缝差异。}

\subsection{目标原生模型的编辑再生检查}
\label{subsec:6-2-3}
\outlinehint{直接在目标系统修改尺寸、参考面或路径，检查相关特征更新；重新转换源模型不能替代此项。}

\subsection{成功率、失败分类与耗时统计}
\label{subsec:6-2-4}
\outlinehint{分别统计支持范围、有效执行和验证通过；明确分母，区分未支持与失败、无法比较与确定不同，并分阶段记录耗时。}

\section{五类特征的基本转换实验}
\label{sec:6-3}

\subsection{草图与基准面转换实验}
\label{subsec:6-3-1}
\outlinehint{选取倾斜放置、圆弧方向、约束位点、偏置及角度基准面，展示参数、特征树和必要几何比较。}

\subsection{拉伸与旋转转换实验}
\label{subsec:6-3-2}
\outlinehint{覆盖单向、双向、对称、参考终止、薄壁和反向轴等场景，分别核对实际范围、材料位置和引用关系。}

\subsection{扫掠转换实验}
\label{subsec:6-3-3}
\outlinehint{围绕三类截面、非圆轮廓、路径反向和圆管内孔设计样例，区分公共表达范围与各路线实际可执行范围。}

\section{组合模型与编辑能力实验}
\label{sec:6-4}

\subsection{连续依赖模型的转换验证}
\label{subsec:6-4-1}
\outlinehint{使用基准面、草图、成形特征和后续引用构成的连续模型，检查引用恢复及误差传播，列出源目标节点对应。}

\subsection{尺寸与上游引用修改后的再生验证}
\label{subsec:6-4-2}
\outlinehint{修改拉伸深度、旋转角度、基准偏置或扫掠路径，记录初始一致性与修改后的预期结果，定位关联丢失。}

\section{对比实验与综合分析}
\label{sec:6-5}

\subsection{关键处理方法的对照实验}
\label{subsec:6-5-1}
\outlinehint{对比是否进行方向归一化、实际坐标转换、薄壁位置恢复和几何分段处理，控制其他条件；未执行的对照不得写成结果。}

\subsection{转换效率与模型规模影响}
\label{subsec:6-5-2}
\outlinehint{统计读取、统一映射、创建、再生和比较耗时，说明进程启动、首次运行、样本量与模型复杂度。}

\subsection{失败案例、能力边界与适用性分析}
\label{subsec:6-5-3}
\outlinehint{围绕读写不对称、参数简化、引用歧义与目标功能限制解释失败，给出结论适用的路线、子类型和条件。}

\section{本章小结}
\label{sec:6-6}

\outlinehint{归纳本章能够支持的认识，并说明与下一章的衔接；不重复列举所有接口或扩展未证明结论。}
